В первой строке входного файла содержится одно число $N$ --- количество аминокислот. Далее следует $N$ непустых строк, каждая из которых является описанием аминокислоты. Суммарная длина строк $S$ не превосходит $10^6$.

В следующей строке содержится исходный белок, длина $L$ которого не превосходит $10^6$.

Далее содержится число $M$ --- количество простых запросов ($0 \leqslant M \leqslant 10^5$). Каждый простой запрос описывается двумя числами --- $l$ и $r$, ($1 \leqslant l \leqslant r \leqslant L$), обозначающими искомый фрагмент $[l, r]$, который необходимо разбить.

Далее содержится число $K$ --- количество запросов-серий ($0 \leqslant K \leqslant 10^5$, в некоторых группах тестов запросы-серии отсутствуют и $K=0$). Каждая серия описывается восемью целыми числами $T, A, B, C, D, E, l_1, r_1$ ($1 \leqslant A, B, C, D \leqslant 100$, $1 \leqslant T, E, l_1, r_1 \leqslant 10^7$). Такая серия описывает $T$ простых запросов. Запрос с номером $i$ выглядит как фрагмент $[\min(l_i \bmod L + 1, r_i \bmod L + 1), \max(l_i \bmod L + 1, r_i \bmod L + 1)]$. Числа $l_i$ и $r_i$ для $i \geqslant 2$ определяются следующим образом. Пусть $ans$ обозначает ответ на запрос с номером $i-1$ в этой серии. Тогда $l_i = (A \cdot l_{i-1} + B) \bmod E, r_i = (C \cdot r_{i-1} + D + ans) \bmod E$. Суммарное число запросов в сериях не превосходит $10^7$.

Все строки во входном файле состоят только из маленьких латинских букв. Гарантируется, что ни одна аминокислота не является префиксом другой.